\documentclass[12pt]{article}
\usepackage[margin=0.6in,top=1.15in,bottom=0.55in]{geometry}
\usepackage{lmodern}
\usepackage{microtype}
\usepackage{titlesec}
\usepackage{enumitem}
\usepackage[hidelinks]{hyperref}
\usepackage{../../latex/grader-worksheet}
\setlength{\parindent}{0pt}
\setlength{\parskip}{0.25em}
\titleformat{\section}{\large\bfseries\scshape}{}{0pt}{}
\GraderSetup{assignment-id=U1L18LL,template-revision=1}
\begin{document}
\begin{center}
{\LARGE\bfseries Week 18 Logic Lab}\par\vspace{0.05em}
{\large\scshape Portfolio Revision Workshop}\par\vspace{0.1em}
\rule{0.78\textwidth}{0.5pt}\par\vspace{0.15em}
\end{center}
\noindent\textbf{Purpose:} Build one verifiable revision trail: select an artifact, recover its reasoning, diagnose a precise weakness, revise substantively, and explain the growth shown.
\vspace{0.3em}

\section*{Part 1: Select and Recover}
\textbf{Question 1.1}\\[-0.15em]
Identify one completed portfolio artifact by assignment or title. State its original claim or task and quote or accurately describe one feature of your original reasoning.
\GraderAnswerBox{Part 1 Q1}{2.45in}

\textbf{Question 1.2}\\[-0.15em]
Why is this artifact a useful revision candidate? Name the before-and-after growth it could make visible; do not answer only that the score was low or high.
\GraderAnswerBox{Part 1 Q2}{2.35in}

\newpage
\section*{Part 2: Diagnose with a Criterion}
\textbf{Question 2.1}\\[-0.15em]
Recover the artifact's reasoning: state its conclusion, support, and one hidden assumption or inference connecting them.
\GraderAnswerBox{Part 2 Q1}{2.55in}

\textbf{Question 2.2}\\[-0.15em]
Diagnose one exact weakness or limit. Name the fitting Whately criterion and point to the feature of the original that the criterion tests.
\GraderAnswerBox{Part 2 Q2}{2.55in}

\newpage
\section*{Part 3: Revise the Reasoning}
\textbf{Question 3.1}\\[-0.15em]
Write a revised claim or conclusion that responds to the diagnosis while preserving what the original evidence can honestly support.
\GraderAnswerBox{Part 3 Q1}{2.45in}

\textbf{Question 3.2}\\[-0.15em]
Revise the support or inference. Add, remove, clarify, qualify, or reconnect material so that the revised reasoning earns the preceding claim.
\GraderAnswerBox{Part 3 Q2}{2.55in}

\newpage
\section*{Part 4: Caption and Connect}
\textbf{Question 4.1}\\[-0.15em]
Write a reflective caption for the revision trail. Explain what the original attempted, what changed, which logical habit guided the change, and what the revision now shows about your judgment.
\GraderAnswerBox{Part 4 Q1}{3.15in}

\textbf{Question 4.2}\\[-0.15em]
Connect this artifact to one other portfolio artifact. State one recurring habit or pattern of growth that both pieces can support, and name the evidence in each.
\GraderAnswerBox{Part 4 Q2}{2.3in}

\newpage
\section*{Part 5: Review}
\textbf{Question 5.1}\\[-0.15em]
A student revises ``Three crowded buses prove the entire transit system is unreliable'' to ``These three trips show a recurring crowding problem on one route.'' Use a prior-week tool to explain one logical improvement and one remaining evidence need.
\GraderAnswerBox{Part 5 Q1}{2.2in}

\textbf{Question 5.2}\\[-0.15em]
An essay says, ``The principal's speech was emotional, so none of its safety claims had evidence.'' Use a prior-week rhetoric tool to repair this judgment without treating emotion as proof.
\GraderAnswerBox{Part 5 Q2}{2.2in}
\end{document}
